\section*{Chapter \ref{ch:pl:distributed-compute-and-schedulers} --- Distributed Compute and Schedulers}

\begin{solution}{exo:pl:distributed-compute-and-schedulers:1}
Longest first puts 3 and 3 on the two cores, then 2 and 2, then the last 2 on the first core free: 7. The best is 6 (3 and 3 on one core, 2,
2 and 2 on the other). The ratio $7/6$ equals Graham's bound $4/3 - 1/(3m)$ for $m = 2$: the example is the worst case.
\end{solution}

\begin{solution}{exo:pl:distributed-compute-and-schedulers:2}
Thirty times fifteen minutes, 7.5~hours, on a normal node; four times that, 30~hours, on the slow node.
\end{solution}

\begin{solution}{exo:pl:distributed-compute-and-schedulers:3}
The total work is 3\,649 core-hours; spread over 384 cores that is 9.50~hours, longer than the longest task (7.5~hours on a normal node) and
than any chain (the flat sweep has none). The bound takes the largest of the three.
\end{solution}

\begin{solution}{exo:pl:distributed-compute-and-schedulers:4}
A healthy task runs its true duration, which exceeds 2.5 times its estimate only when the estimate's error $e^{0.3z}$ is below $1/2.5$, that is
$z > 3.05$: a probability of 0.11\%. At 1.5 the threshold is passed whenever $z > 1.35$, about 9\% of tasks: hundreds of useless backups,
each taking a core from the queue and wasting work.
\end{solution}

\begin{solution}{exo:pl:distributed-compute-and-schedulers:5}
\omterm{def:pl:distributed-compute-and-schedulers:stealing}{Work stealing} keeps every core busy while work remains, but it takes tasks in the order they were dealt; the long tasks, submitted last, start
late and end late, as under first in, first out. Balance is not order: only a policy that starts the longest tasks first avoids the tail.
\end{solution}

\begin{solution}{exo:pl:distributed-compute-and-schedulers:6}
In this run it ends earlier, at 16.68~hours instead of 19.79: the second team's 500 tasks are only about 5\% of the work, and the first
team's end is set by where and when its forty long tasks run, which the interleaving changes --- here to its advantage, in another run
perhaps not.
\end{solution}

\begin{solution}{exo:pl:distributed-compute-and-schedulers:7}
First in, first out improves from 19.79 to 14.41~hours (13.69 with backups): the long tasks start at once. It still trails longest first,
which is unchanged at 11.80 (10.02 with backups) because it orders by estimate, not by submission: the other long-tailed tasks are still in
random order under first in, first out.
\end{solution}

\begin{solution}{exo:pl:distributed-compute-and-schedulers:8}
Utilisation over a night says nothing about how the capacity was used: the sweep held 475.0 node-hours for 3\,649 core-hours of work because of
the order of its tasks, and it was almost idle for three and a half hours. On half the cluster the lower bound doubles to 19~hours; the fix is the
schedule (longest first, backups), which ends near the bound on the cluster as it is.
\end{solution}

\begin{solution}{pb:pl:distributed-compute-and-schedulers:1}
\begin{enumerate}
\item Each backtest reads its inputs and writes its result without talking to the others.
\item From 24 small backtests of chapter~11 run on this laptop through the local executor: the spread of log durations, 0.76.
\item 3\,649 core-hours; 9.50~hours on 384 cores.
\item At the end, because a grid over parameters submits its largest settings last.
\item Intrinsic length, a slow node, and failure with a retry from the start.
\item At most $2 - 1/m$ times the best \omterm{def:pl:distributed-compute-and-schedulers:makespan}{makespan} for any order; at most $4/3 - 1/(3m)$ for longest first.
\item Estimates of duration.
\item It balances load but keeps the dealt order, so the long tasks still start last.
\item It serves the team using less than its share first, so a small late sweep is not queued behind a large one.
\item When a task is slow because of where it runs (a slow node); it wastes when estimates are poor and the threshold fires on healthy tasks.
\item First in, first out 19.79 (20.32 with backups); longest first 11.80 (10.02); \omterm{def:pl:distributed-compute-and-schedulers:stealing}{work stealing} 16.33 (16.40); longest first without the slow
node 10.32.
\item About three and a half hours: fewer than a tenth of the cores were busy from 16.3~hours to 19.8.
\item 5.5\% (10.02 against 9.50), because of the slow node and the estimation error; without the slow node but without backups, 10.32.
\item Half its tasks are done 0.59~hours after submission instead of 8.28.
\item The graph computes the features once (3\,514 instead of 6\,922 core-hours, 11.4 instead of 19.3~hours); the cache saves the hundred
feature stages in the next sweep (3\,475 core-hours, 10.3~hours).
\item \textbf{Named result.} On 384 cores the 10\,000-task sweep takes 19.79~hours and 475.0 node-hours first in, first out, 11.80 and 283.3
longest first, 16.33 and 391.9 with \omterm{def:pl:distributed-compute-and-schedulers:stealing}{work stealing}; backups after 2.5 times the estimate save longest first 1.8~hours and cost first in,
first out 0.5; the best schedule, longest first with backups, ends 5.5\% above the lower bound of 9.50~hours.
\item Longest first on estimates by default, backups with a threshold set from the estimates' error, fair share between teams.
\item An estimate of each task's duration (or the parameters that predict it), and cache keys for shared stages.
\item From the lower bound of the typical sweep and the deadline: cores = work / (deadline $-$ longest task margin), then check with the
simulator under the policy actually used.
\item Its last tasks: their length, where they run and when they start.
\end{enumerate}
\end{solution}

\begin{solution}{iq:pl:distributed-compute-and-schedulers:1}
Find out why they are long: intrinsic (start them first with longest-first ordering on estimates), a slow machine (backups on another node),
or failures (checkpoints, retries); then check the \omterm{def:pl:distributed-compute-and-schedulers:makespan}{makespan} against the lower bound.
\iqlookfor{Diagnose the tail before adding capacity.}
\end{solution}

\begin{solution}{iq:pl:distributed-compute-and-schedulers:2}
Running a second copy of a task that is late, on another machine, keeping the first to finish. It is a bad idea when tasks are late because
they are long rather than because of where they run, when estimates are too poor to tell, or when the work has side effects that two copies
would duplicate.
\iqlookfor{Thresholds from estimate error; idempotent tasks.}
\end{solution}

\begin{solution}{iq:pl:distributed-compute-and-schedulers:3}
Fair share with targets per team, so that a team under its share is served first; limits on concurrent cores per user; and priority for
short interactive work over long sweeps.
\iqlookfor{Shares, not first come first served.}
\end{solution}

\begin{solution}{iq:pl:distributed-compute-and-schedulers:4}
As a \omterm{def:pl:distributed-compute-and-schedulers:graph}{task graph}: the feature computations as shared stages with content-addressed cache keys, the backtests depending on them; then schedule
the graph, longest estimated tasks first.
\iqlookfor{Graph, cache keys, reuse across sweeps.}
\end{solution}

\begin{solution}{iq:pl:distributed-compute-and-schedulers:5}
\omterm{def:pl:distributed-compute-and-schedulers:graph}{Task graphs} with resource requests and estimates; a queue with fair share between teams and longest-first ordering within a sweep; retries
with checkpoints; backups with a calibrated threshold; a content-addressed cache; accounting of node-hours per team; and a simulator of the
cluster to test policies before they are deployed.
\iqlookfor{Policy chosen by measurement and simulation, not by default.}
\end{solution}
